\chapter{\texorpdfstring{$L^p$}{Lp} 空间}\label{ch:b3:lp}

勒贝格积分是为分析而建立的；$L^p$ 空间
正是这种分析的用武之地。它们是巴拿赫空间
（里斯–费舍尔） --- 正是
\cref{exo:b3:complete:1}所示连续函数所缺的
那种完备化 --- 同时还支持一种光滑化技术，
\emph{与磨光核的卷积}，它可将任意
$L^p$ 函数用 $\mathcal C^\infty$ 函数逼近。本章证明积分形式的
赫尔德不等式与闵可夫斯基不等式、完备性、
稠密性定理，以及正则化方法，最后
以 $L^p$ 族的包含与插值关系作结。
本章始终设 $(X, \mathcal A, \mu)$ 为测度空间，
函数取复值；在 $\R^d$ 上，测度为
$\lambda_d$。

\section{定义；赫尔德与闵可夫斯基}

\begin{definition}\label{def:b3:lp:space}
对 $1 \leq p < \infty$，$\mathcal L^p(\mu)$ 是所有
可测函数 $f$ 满足 $\norm f_p = \bigl(\int\abs
f^p\dd\mu\bigr)^{1/p} < \infty$ 的集合，而 $\mathcal L^\infty(\mu)$
则为在某零测集外有界的 $f$ 构成的集合，并记其范数 $\norm
f_\infty$ 为\emph{本性上确界} --- 即最小的 $M$，使得
$\abs f \leq M$ 几乎处处成立（下确界可以取到：取
$M + \frac1n$ 对应的零测集之交）。由于 $\norm f_p = 0$ 只能推出 $f
= 0$ \emph{几乎处处}（见\cref{exo:b3:lebesgue:5}），于是定义
\[
L^p(\mu) = \mathcal L^p(\mu)/\{f = 0 \text{ 几乎处处}\} :
\]
其元素是函数模去零测集所得的等价类，而
$\norm\cdot_p$ 是 $L^p$ 上
的真正范数。\index{$L^p$ 空间}
\end{definition}

\begin{theorem}[赫尔德不等式]\label{thm:b3:lp:holder}
设 $1 \leq p, q \leq \infty$，且 $\frac1p + \frac1q = 1$
（共轭指数）。对可测函数 $f, g$：\index{赫尔德
不等式}
\[
\norm{fg}_1 \leq \norm f_p\,\norm g_q ,
\]
等号成立（当 $1 < p < \infty$、范数有限、$f,g\ne0$ 时）当且仅当
$\abs f^p$ 与 $\abs g^q$ 几乎处处成比例。
\end{theorem}

\begin{proof}
当 $\{p, q\} = \{1, \infty\}$ 时结论直接可得（$\abs{fg} \leq
\norm g_\infty\abs f$ 几乎处处成立）。设 $1 < p < \infty$，
不妨规范化为 $\norm f_p = \norm g_q = 1$（由齐次性；零或无穷
范数的情形平凡）。杨不等式 $ab \leq \frac{a^p}p +
\frac{b^q}q$（$a, b \geq 0$；由 $\ln$ 的凹性，参见
\cref{pb:b3:banach:1}）得逐点不等式 $\abs{f g} \leq
\frac{\abs f^p}p + \frac{\abs g^q}q$；两边积分，得 $\norm{fg}_1
\leq \frac1p + \frac1q = 1$。等号成立迫使杨不等式几乎处处取等，
即 $\abs f^p = \abs g^q$ 几乎处处成立（在规范化之后；
撤销规范化后即为成比例）。
\end{proof}

\begin{theorem}[闵可夫斯基不等式]\label{thm:b3:lp:minkowski}
对 $1 \leq p \leq \infty$，有 $\norm{f + g}_p \leq \norm f_p +
\norm g_p$。\index{闵可夫斯基不等式}
\end{theorem}

\begin{proof}
$p = 1, \infty$：逐点或几乎处处的三角不等式。设 $1 <
p < \infty$，假设 $\norm{f+g}_p < \infty$（否则可用
$\abs{f+g}^p \leq 2^{p-1}(\abs f^p + \abs g^p)$，由
$t^p$ 的凸性，可知右边有限时左边亦有限）。
于是
\[
\norm{f{+}g}_p^p \leq \int\abs f\,\abs{f{+}g}^{p-1} +
\int\abs g\,\abs{f{+}g}^{p-1}
\leq \bigl(\norm f_p + \norm g_p\bigr)\,
\bigl\|\abs{f{+}g}^{p-1}\bigr\|_q
\]
由赫尔德不等式，且 $\norm{\abs{f+g}^{p-1}}_q =
\norm{f+g}_p^{p/q}$，因为 $(p-1)q = p$；两边同除以
$\norm{f+g}_p^{p/q}$（若非零，否则平凡），并利用 $p -
\frac pq = 1$。
\end{proof}

\section{完备性及相关结论}

\begin{theorem}[里斯–费舍尔]\label{thm:b3:lp:complete}
对 $1 \leq p \leq \infty$，$L^p(\mu)$ 是巴拿赫空间。
此外，$L^p$ 中每个收敛的序列都有一个子列，
\emph{几乎处处}收敛（附以 $L^p$ 控制函数，
仅当 $p < \infty$ 时）。\index{里斯–费舍尔定理}
\end{theorem}

\begin{proof}
$p = \infty$：$\norm\cdot_\infty$ 意义下的柯西序列在
除去单个零测集（可数个零测集的并）外一致柯西：
它在该集外一致收敛，证毕。设 $p < \infty$。由
\cref{exo:b3:complete:1}(b)，只需证明绝对收敛的级数收敛：
设 $\sum\norm{f_k}_p = M < \infty$。令
$G_n = \sum_{k \leq n}\abs{f_k}$ 与 $G = \sum_k\abs{f_k}$
（逐点在 $[0,\infty]$ 中取值）：由闵可夫斯基不等式，$\norm{G_n}_p \leq
M$，由单调收敛定理（$G_n^p \nearrow G^p$）得 $\int G^p \leq M^p$：
$G < \infty$ 几乎处处成立，故级数 $\sum f_k(x)$
在几乎处处的 $x$ 处绝对收敛；记其和为 $S(x)$
（在零测集上任取值）。于是 $\abs{S - \sum_{k\leq n}f_k}^p \leq (2G)^p
\in L^1$，由控制收敛定理得 $\norm{S - \sum_{k\leq n}f_k}_p \to 0$：
即级数在 $L^p$ 中收敛。

关于子列的论断：若 $f_n \to f$ 于 $L^p$，取 $n_k$
使 $\norm{f_{n_{k+1}} - f_{n_k}}_p \leq 2^{-k}$，考虑级数
$\sum(f_{n_{k+1}} - f_{n_k})$ 可归入前面的论证，
即它几乎处处绝对收敛，并以某个 $G \in
L^p$ 为控制，故 $f_{n_k} \to f_{n_1} + \sum(\cdots) $ 几乎处处成立，且这一
几乎处处极限必为 $f$ 的代表元（二者均为 $L^p$
极限）。控制函数为 $\abs{f_{n_k}} \leq \abs{f_{n_1}} + G$。
\end{proof}

\begin{remark}\label{rem:b3:lp:modes}
$L^p$ 收敛不蕴含几乎处处收敛（
即\emph{打字机序列}，见\cref{exo:b3:lp:3}），反之亦不成立
（逃逸凸起）：两种收敛仅
通过子列与控制相联系。牢记
\cref{exo:b3:lp:3}中的反例，是最好的“疫苗”。
\end{remark}

\section{稠密性定理}

\begin{theorem}\label{thm:b3:lp:density}
设 $1 \leq p < \infty$。
\begin{enumerate}
  \item 简单函数（支集具有有限测度）在
        $L^p(\mu)$ 中稠密。
  \item 在 $L^p(\R^d)$ 中，连续紧支撑函数
        类 $\mathcal C_c(\R^d)$ 稠密。
  \item 平移在 $L^p(\R^d)$ 上连续：记
        $\tau_hf = f(\cdot - h)$，$\norm{\tau_hf - f}_p \to 0$
        当 $h \to 0$ 时。
\end{enumerate}
以上三条对 $p = \infty$ 均不成立。
\end{theorem}

\begin{proof}
(1) 设 $f \geq 0$，由
\cref{thm:b3:lebesgue:approximation}中的二进列 $s_n \nearrow f$ 满足 $\abs{f - s_n}^p
\leq f^p \in L^1$，由控制收敛定理。（每个 $s_n \leq f$ 属于 $L^p$，
其水平集在取值为正处具有有限测度：
$\mu(s_n \geq c) \leq c^{-p}\int f^p$。）将
一般的 $f$ 分解为四个非负部分。

(2) 由 (1)，只需逼近 $\mathbf 1_A$，其中 $A$ 为博雷尔集，
且 $\lambda_d(A) < \infty$。由正则性（证明同
\cref{thm:b3:measure:regularity}）得紧集 $K \subseteq A
\subseteq U$ 为开集，使 $\lambda_d(U\setminus K) <
\varepsilon$；取乌雷松函数
\[
\varphi(x) = \frac{d(x, \R^d\setminus U)}{d(x, \R^d\setminus U)
+ d(x, K)}
\]
该函数连续，$1$ 在 $K$ 上，$0$ 在 $U$ 外，
且可取为紧支撑的（先将 $U$ 缩小为有界开集）。于是有
$\norm{\mathbf 1_A - \varphi}_p^p \leq \lambda_d(U\setminus K)
< \varepsilon$。

(3) 对 $g \in \mathcal C_c$，由一致连续性得
$\norm{\tau_hg - g}_\infty \to 0$，且支集在
当 $\abs h \leq 1$ 时位于一固定紧集内，故 $\norm{\tau_hg - g}_p \to 0$。对
一般的 $f$，取 $g \in \mathcal C_c$，使 $\norm{f - g}_p <
\varepsilon$；则 $\norm{\tau_hf - f}_p \leq 2\norm{f - g}_p +
\norm{\tau_hg - g}_p$（范数的平移不变性）。

对 $p = \infty$，用连续函数一致逼近 $\mathbf
1_{\intoo0\infty}$ 不可能
（存在跳跃），且 $\norm{\tau_h\mathbf 1_{\intoo0\infty} - \mathbf
1_{\intoo0\infty}}_\infty = 1$，当 $h \neq 0$ 时。
\end{proof}

\section{卷积与正则化}

\begin{theorem}[杨不等式]\label{thm:b3:lp:young}
设 $1 \leq p \leq \infty$，$f \in L^1(\R^d)$，$g \in
L^p(\R^d)$。则 $f * g$ 几乎处处有定义，属于 $L^p$，
且\index{杨不等式}
\[
\norm{f * g}_p \leq \norm f_1\,\norm g_p .
\]
\end{theorem}

\begin{proof}
$p = \infty$：直接估计。$p = 1$：
由\cref{thm:b3:product:convolution}。设 $1 < p < \infty$，$q$
为其共轭指数。将 $\abs{f(y)} = \abs{f(y)}^{1/q}\cdot
\abs{f(y)}^{1/p}$，并应用赫尔德不等式：
\[
\int\abs{f(y)}\,\abs{g(x{-}y)}\,\dd y
\leq \Bigl(\int\abs f\Bigr)^{1/q}
\Bigl(\int\abs{f(y)}\,\abs{g(x - y)}^p\,\dd y\Bigr)^{1/p} .
\]
两边取 $p$ 次方并对 $x$ 积分；对第二个因子应用托内利定理，
得 $\norm f_1^{p/q}\cdot\norm f_1\norm g_p^p$，
即 $\norm{f*g}_p^p \leq \norm f_1^{1 + p/q}\norm g_p^p =
(\norm f_1\norm g_p)^p$ --- 而托内利积分的有限性
保证几乎处处绝对收敛，与
\cref{thm:b3:product:convolution}中的论证相同。
\end{proof}

\begin{definition}[磨光子]\label{def:b3:lp:mollifier}
函数\index{磨光子}\index{凸起函数}
\[
\rho(x) = \begin{cases}
c\,\exp\Bigl(-\dfrac{1}{1 - \norm x^2}\Bigr) & \norm x < 1,\\
0 & \norm x \geq 1,
\end{cases}
\]
其中常数 $c$ 使 $\int\rho = 1$，它是 $\mathcal C^\infty$ 函数
（定义于 $\R^d$ 上）：关键在于 $t \mapsto \eu^{-1/t}\mathbf
1_{t>0}$ 为 $\mathcal C^\infty$ 函数（在 $\R$ 上），其各阶导数在
$0^+$ 处均为 $0$（每个导数形如 $P(1/t)\eu^{-1/t}$，
$P$ 为多项式，它趋于 $0$；对阶数用归纳法）。对
$\varepsilon > 0$，令 $\rho_\varepsilon(x) =
\varepsilon^{-d}\rho(x/\varepsilon)$，其支集含于 $\bar B(0,
\varepsilon)$ 内，积分仍为 $1$。
\end{definition}

\begin{theorem}[正则化]\label{thm:b3:lp:regularization}
设 $1 \leq p < \infty$，$f \in L^p(\R^d)$。则：
\begin{enumerate}
  \item $f * \rho_\varepsilon \in \mathcal C^\infty(\R^d)$，
        且 $\partial^\alpha(f * \rho_\varepsilon) = f *
        \partial^\alpha\rho_\varepsilon$；
  \item $\norm{f * \rho_\varepsilon - f}_p \to 0$，当
        $\varepsilon \to 0$ 时；
  \item 因此 $\mathcal C^\infty_c(\R^d)$ 在
        $L^p(\R^d)$ 中稠密。
\end{enumerate}
\end{theorem}

\begin{proof}
(1) 积分号下求导
（\cref{thm:b3:lebesgue:paramdiff}）对 $x$ 求导：设 $x$ 位于球 
$B$ 内，$\abs{\partial_{x_i}\rho_\varepsilon(x - y)} \leq
C_\varepsilon\,\mathbf 1_{K}(y)$，其中 $K$ 为紧集（$y$
在 $\varepsilon$ 之内、距 $B$ 不远），且 $\abs f\,\mathbf 1_K \in L^1$
（对 $\mathbf 1_K$ 用赫尔德不等式）：定理适用；
对高阶导数逐次应用。

(2) 由 $\int\rho_\varepsilon = 1$，
\[
(f * \rho_\varepsilon)(x) - f(x)
= \int \bigl(f(x - y) - f(x)\bigr)\rho_\varepsilon(y)\,\dd y ,
\]
由闵可夫斯基\emph{积分}不等式 --- 或直接用
赫尔德或延森不等式，配以概率测度
$\rho_\varepsilon\dd y$，并用托内利定理 ---
\[
\norm{f*\rho_\varepsilon - f}_p^p
\leq \int\Bigl(\int\abs{f(x-y) -
f(x)}^p\dd x\Bigr)\rho_\varepsilon(y)\,\dd y
= \int \norm{\tau_yf - f}_p^p\;\rho_\varepsilon(y)\,\dd y
\]
（中间一步：应用延森不等式，
见\cref{exo:b3:lp:10}，对内层 $y$ 积分，再用托内利定理）。
被积函数支撑于 $\norm y \leq \varepsilon$ 内，并且
在其上一致地趋于 $0$（$\varepsilon \to 0$ 时），
由\cref{thm:b3:lp:density}(3)可知，整个式子趋于
$0$。

(3) 对 $f$，取 $g \in \mathcal C_c$ 逼近它
（由\cref{thm:b3:lp:density}(2)），再取 $g$ 的磨光 $g *
\rho_\varepsilon \in \mathcal C_c^\infty$（紧支集：
两支集之和）。
\end{proof}

\begin{example}[磨光 $\abs x$，附速率]
\label{ex:b3:lp:mollifyexample}
取 $f(x) = \abs x$，定义于 $\R$ 上（局部 $L^1$；定理在每个有界窗口上
均适用），以及一个对称的磨光子
$\rho_\varepsilon$，则
\[
f_\varepsilon(x) = (f * \rho_\varepsilon)(x)
= \int\abs{x - y}\,\rho_\varepsilon(y)\,\dd y
\]
是 $\mathcal C^\infty$ 函数；远离尖点时，什么也不发生：
当 $\abs x \geq \varepsilon$ 时，$\abs{x - y}$ 关于 $x$ 为线性
在 $\rho_\varepsilon$ 的支集上，从而 $f_\varepsilon(x) =
\abs x$ \emph{恰好}（对称性消去了修正项）。
在 $0$ 附近，磨光的代价恰为
\[
0 \leq f_\varepsilon(0) = \int\abs
y\,\rho_\varepsilon(y)\,\dd y \leq \varepsilon,
\qquad
\norm{f_\varepsilon - f}_\infty \leq \varepsilon :
\]
逼近误差局限于奇点的
$\varepsilon$ 邻域内，且与奇点的
大小同阶。同时 $f_\varepsilon'' \geq 0$ 处处成立
（$f$ 为凸函数，与 $\rho_\varepsilon
\geq 0$ 的卷积保持凸性），且 $\int f_\varepsilon'' =
f_\varepsilon'(\infty) - f_\varepsilon'(-\infty) = 2$：二阶导数
是一个质量为 $2$ 的凸块，被压缩进宽度
$O(\varepsilon)$ 内，因此 $\norm{f_\varepsilon''}_\infty \gtrsim
\varepsilon^{-1}$。磨光是一种权衡：一致误差
$O(\varepsilon)$ 对应导数的爆破
$O(\varepsilon^{-1})$——这正是
量化分析（插值不等式，
\cref{pb:b3:lp:1}的思路）所形式化的精确交换率。
\end{example}

\begin{corollary}[变分法基本
引理]\label{cor:b3:lp:fundlemma}
设 $f \in L^1_{\mathrm{loc}}(\R^d)$，在紧集上可积
且 $\int f\varphi = 0$，对每个 $\varphi \in \mathcal
C^\infty_c(\R^d)$。则 $f = 0$ 几乎处处\index{变分法基本
引理}
\end{corollary}

\begin{proof}
固定球 $B = B(0, R)$，令 $g = f\mathbf 1_{B(0, R+1)}
\in L^1$。对 $x \in B$ 及 $\varepsilon < 1$，有 $(g *
\rho_\varepsilon)(x) = \int f(y)\rho_\varepsilon(x - y)\dd y =
0$，检验函数为 $y \mapsto \rho_\varepsilon(x-y)
\in \mathcal C_c^\infty$。但 $g * \rho_\varepsilon \to g$ 在
$L^1$ 中收敛（由\cref{thm:b3:lp:regularization}）：$g = 0$ 在 $B$ 上几乎处处成立；
由这些球穷竭 $\R^d$。
\end{proof}

\section{\texorpdfstring{$L^p$}{Lp}的版图}

\begin{proposition}\label{prop:b3:lp:inclusions}
(a) 若 $\mu(X) < \infty$，且 $1 \leq p \leq q \leq \infty$，
则 $L^q \subseteq L^p$，且 $\norm f_p \leq
\mu(X)^{\frac1p - \frac1q}\,\norm f_q$。
(b) 在 $\R^d$ 上（测度无限），不存在包含关系：
即当 $p \neq q$ 时，存在属于 $L^p\setminus L^q$ 的函数。
(c)（插值）若 $p < r < q$，且 $\alpha \in \intoo01$
由 $\frac1r = \frac\alpha p + \frac{1 - \alpha}q$ 定义，
则
\[
\norm f_r \leq \norm f_p^{\alpha}\,\norm f_q^{1 - \alpha} ;
\]
特别地，$L^p \cap L^q \subseteq L^r$。
\end{proposition}

\begin{proof}
(a) 用指数为 $\frac qp$ 及其共轭的赫尔德不等式：
$\int\abs f^p\cdot 1 \leq \norm{\abs f^p}_{q/p}\,\norm
1_{(q/p)'} = \norm f_q^p\,\mu(X)^{1 - p/q}$（$q = \infty$ 时
直接可得）。(b) 在 $0$ 附近及 $\infty$ 附近，幂 $x^{-\alpha}$
起校准作用，见\cref{exo:b3:lp:2}。(c) 写 $\abs f^r = \abs
f^{r\alpha}\abs f^{r(1-\alpha)}$，并用赫尔德不等式，取
共轭对 $\frac p{r\alpha}$、$\frac q{r(1-\alpha)}$，
（由 $\alpha$ 的定义，二者恰好共轭）：
$\int\abs f^r \leq \norm f_p^{r\alpha}\norm f_q^{r(1 -
\alpha)}$。
\end{proof}

\begin{method}\label{met:b3:lp:toolkit}
下文统一使用的 $L^p$ 工具如下：欲对一切
$f \in L^p$ 证明某恒等式或不等式，先在稠密类上证明，
稠密类取 $\mathcal C_c^\infty$（见
\cref{thm:b3:lp:regularization}），再由连续性延拓
（\cref{thm:b3:complete:extension}，两端都是
$L^p$-连续的）；要证 $f = 0$，用 $\mathcal
C_c^\infty$ 检验（\cref{cor:b3:lp:fundlemma}）；要获得光滑性，
则作卷积；要互换指数，用赫尔德不等式与插值。
在\cref{ch:b3:fouriertransform}中的傅里叶理论，
整体上正是本方法的一个长篇应用。
\end{method}

\section{练习}

\begin{exercise}[$\star$]\label{exo:b3:lp:1}
(a) 叙述并证明 $L^2(\mu)$ 中的柯西–施瓦茨不等式，
它是赫尔德不等式在 $p = q = 2$ 时的情形。
在\emph{概率空间}上，证明映射 $p \mapsto \norm f_p$
是非减的。
(c) 当 $p = 1$、$q = \infty$ 时，赫尔德不等式何时取等号？
\end{exercise}

\begin{exercise}[$\star$]\label{exo:b3:lp:2}
对于哪些 $p \in \intco1\infty$，下列函数属于
$L^p$？
\[
x^{-1/2}\mathbf 1_{\intoo01},\qquad
x^{-1/2}\mathbf 1_{\intoo1\infty},\qquad
\frac{1}{x^{1/2}(1 + \abs{\ln x})}\ \text{ 在 } \intoo01,
\qquad
\frac1{1 + \abs x}\ \text{ 在 } \R .
\]
由此得出：在 $\intoo01$ 上，较小的 $p$ 更易属于；在
$\intoo1\infty$ 上，较大的 $p$ 更易属于；且任何 $L^p$ 都
不包含另一个（在 $\R$ 上）。
\end{exercise}

\begin{exercise}[$\star\star$]\label{exo:b3:lp:3}
（打字机）列举二进区间 $I_1 =
\intcc01$, $I_2 = \intcc0{\frac12}$, $I_3 = \intcc{\frac12}1$,
$I_4 = \intcc0{\frac14}$, \dots\ 并令 $f_n = \mathbf
1_{I_n}$.
(a) 证明 $f_n \to 0$ 在每个 $L^p(\intcc01)$ 中成立（$p < \infty$），
但 $(f_n(x))$ 对\emph{每个} $x \in \intcc01$ 都发散。
(b) 给出
\cref{thm:b3:lp:complete}所保证的几乎处处收敛的子列。
(c) 反之，给出一个几乎处处收敛、
但在 $L^1$ 中不收敛的数列，以及一个在 $L^1$ 中收敛、但在任何 $L^p$（$p > 1$）中都不收敛的数列。
\end{exercise}

\begin{exercise}[$\star\star$]\label{exo:b3:lp:4}
设 $\mu(X) < \infty$，$f \in L^\infty(\mu)$，$f \neq 0$。
证明 $\norm f_p \to \norm f_\infty$（$p \to \infty$）。
\emph{（上界由\cref{prop:b3:lp:inclusions}的 (a) 得出；
下界通过对集合 $\{\abs f > \norm f_\infty -
\varepsilon\}$ 上积分得到，该集合测度为正。）}
\end{exercise}

\begin{exercise}[$\star\star$]\label{exo:b3:lp:5}
(a) 究竟在
\cref{thm:b3:lp:density}(3)的证明中何处用到了 $p < \infty$？
(b) 证明 $f \in L^\infty(\R)$ 满足 $\norm{\tau_hf -
f}_\infty \to 0$ 当且仅当 $f$ 有一个一致连续的
代表元。
\end{exercise}

\begin{exercise}[$\star\star$]\label{exo:b3:lp:6}
设 $p, q$ 为共轭指数，$f \in L^p(\R^d)$，$g \in
L^q(\R^d)$。证明 $f * g$ \emph{处处}有定义，
有界，且 $\norm{f*g}_\infty \leq \norm f_p\norm g_q$，并且
\emph{一致连续}。\emph{（平移的连续性
在 $L^p$ 中；对 $p \in \{1, \infty\}$ 另作处理 --- 当 $p =
\infty$ 时，对 $L^1$ 因子使用平移的连续性。）}
\end{exercise}

\begin{exercise}[$\star\star$]\label{exo:b3:lp:7}
设 $f \in L^1_{\mathrm{loc}}(\intoo ab)$，且
$\int f\varphi' = 0$ 对每个 $\varphi \in \mathcal
C^\infty_c(\intoo ab)$ 成立。证明 $f$ 几乎处处等于某个
常数。\emph{（取定 $\chi \in \mathcal C_c^\infty$，使得
$\int\chi = 1$；任何 $\psi \in \mathcal C_c^\infty$ 满足
$\int\psi = 0$，则它是某个 $\varphi'$；一般检验函数写作
$\psi + (\int\psi)\chi$，并应用
\cref{cor:b3:lp:fundlemma}于 $f - c$，取 $c = \int
f\chi$。）}
\end{exercise}

\begin{exercise}[$\star\star\star$]\label{exo:b3:lp:8}
（光滑乌雷松引理）设 $K \subseteq U \subseteq \R^d$，$K$
为紧集，$U$ 为开集。构造 $\varphi \in \mathcal
C^\infty_c(\R^d)$，满足 $0 \leq \varphi \leq 1$，$\varphi = 1$
在 $K$ 上，$\operatorname{supp}\varphi \subseteq U$。
\emph{（对 $\delta$ 邻域
$K_\delta$（$K$ 的）示性函数，用 $\rho_{\delta/2}$ 光滑化，$\delta$ 充分小。）}
由此推出一个 $\mathcal C^\infty$ 单位分解命题，
适用于被有限个开集覆盖的紧集。
\end{exercise}

\begin{exercise}[$\star\star$]\label{exo:b3:lp:9}
利用插值（\cref{prop:b3:lp:inclusions}(c)）：
(a) 证明 $L^1(\R)\cap L^\infty(\R) \subseteq L^p(\R)$ 对
所有 $p$ 成立，且 $\norm f_p \leq \norm f_1^{1/p}\norm
f_\infty^{1 - 1/p}$；
(b) 证明 $f \mapsto \norm f_p$，对固定的 $f$，
关于 $\frac1p$ 对数凸，并给出一个例子，其中 $f \in L^p$
恰好当 $p$ 位于给定区间 $(p_0, p_1)$ 内时成立。
\end{exercise}

\begin{exercise}[$\star\star$]\label{exo:b3:lp:10}
（延森）设 $\mu$ 为\emph{概率}测度，$f \in
L^1(\mu)$ 为实值函数，且 $\Phi \colon \R \to \R$ 为凸函数。证明
\[
\Phi\Bigl(\int f\,\dd\mu\Bigr) \leq \int \Phi\circ
f\,\dd\mu
\]
\emph{（$\Phi$ 在点 $m = \int f$ 处的支撑线）}。
由此推出算术–几何平均不等式，以及
$p \mapsto \norm f_p$ 的单调性（见
\cref{exo:b3:lp:1}(b)）。\index{延森不等式}
\end{exercise}

\begin{exercise}[$\star\star\star$]\label{exo:b3:lp:11}
（杨氏卷积不等式）设 $1 \leq p, q, r \leq
\infty$，且 $\frac1p + \frac1q = 1 + \frac1r$，$f \in
L^p(\R^d)$，$g \in L^q(\R^d)$。
(a) 证明 $\norm{f * g}_r \leq \norm f_p\,\norm g_q$。
\emph{（由 $p, q,
r$,
\[
\abs{f(y)g(x-y)} = \bigl(\abs f^p\abs g^q\bigr)^{1/r}
\cdot\abs f^{\,p(1/p - 1/r)}\cdot\abs g^{\,q(1/q - 1/r)},
\]
确定的共轭指数，写出上式，并应用三因子赫尔德不等式，指数分别为
$r$、$\frac{pr}{r - p}$、$\frac{qr}{r - q}$；然后
对 $x$ 用托内利定理积分。）}
(b) 验证已知的三个特殊情形：$r =
\infty$（赫尔德不等式，见\cref{exo:b3:lp:6}）；$q = 1$
（可积核卷积的 $L^p$ 稳定性）；
$p = q = 1$（$L^1$ 是卷积代数，见
\cref{thm:b3:product:convolution}).
(c) 为什么不存在 $\frac1p + \frac1q < 1 +
\frac1r$ 的不等式？\emph{（以伸缩 $f_\lambda(x) =
f(\lambda x)$ 作检验，比较两边的标度。）}
\end{exercise}

\begin{exercise}[$\star\star$]\label{exo:b3:lp:12}
（等号成立的情形）(a) 在赫尔德不等式
$\int\abs{fg} \leq \norm f_p\norm g_q$（$1 < p < \infty$）中，
证明等号成立当且仅当 $\abs f^p$ 与 $\abs g^q$
几乎处处成比例。\emph{（考察杨氏
不等式 $ab \leq \frac{a^p}p + \frac{b^q}q$ 的等号情形，即 $a^p
= b^q$.)}
(b) 在闵可夫斯基不等式 $\norm{f + g}_p \leq \norm f_p
+ \norm g_p$（$1 < p < \infty$）中，证明当 $f,
g \neq 0$ 时等号成立迫使 $g = cf$ 几乎处处，且 $c > 0$。
(c) 与 $p = 1$ 和 $p = \infty$ 的情形对比，通过例子描述
那里（范围大得多的）等号成立情形。
\end{exercise}

\section{周末问题：哈代不等式}

\begin{problem}[{周末问题 --- 哈代不等式及其
精确常数}]\label{pb:b3:lp:1}
设 $f \in L^p(\intoo0{+\infty})$，$1 < p < \infty$，定义
\emph{哈代算子}
\[
(Hf)(x) = \frac1x\int_0^x f(t)\,\dd t .
\]
哈代不等式（1920 年）成立如下：\index{哈代不等式}
\[
\norm{Hf}_p \;\leq\; \frac{p}{p-1}\,\norm f_p ,
\]
而常数 $\frac p{p-1}$ 是最优的，且不能取到。
本题先完整证明上述结论，再推广到级数。

\medskip
\noindent\textbf{第一部分 --- 不等式。}首先设 $f
\geq 0$ 的连续函数，在
$\intoo0{+\infty}$ 中具紧支集，并令 $F(x) = \int_0^xf$。
\begin{enumerate}
  \item 证明 $Hf \in L^p$：在 $0$ 附近，$F$ 在
        $0$ 的某一邻域上恒为零；在 $\infty$ 附近 $F$ 有界，从而
        $(Hf)(x) = O(1/x)$，且 $x \mapsto \frac1x$ 属于
        $L^p(\intoo1{+\infty})$，只要 $p > 1$。
  \item 分部积分证明
        \[
        \int_0^\infty \Bigl(\frac Fx\Bigr)^p\dd x
        = \frac{p}{p-1}\int_0^\infty\Bigl(\frac
        Fx\Bigr)^{p-1}f(x)\,\dd x .
        \]
        \emph{（对 $x^{1-p}F^p$ 求导；边界项
        为零 --- 两端都需论证。）}
  \item 对右端应用赫尔德不等式，并推出
        $\norm{Hf}_p \leq \frac p{p-1}\norm f_p$ 对此类 $f$ 成立。
  \item 推广到全体 $L^p$：对 $f \geq 0$，构造
        连续函数 $f_n$，在
        $\intoo0{+\infty}$ 中具紧支集，并使 $0 \leq f_n \nearrow f$ 几乎处处成立
        \emph{（先截断，再单调逼近 ---
        说明构造的合理性）}；然后 $Hf_n \nearrow Hf$
        逐点成立（在平均号内用单调收敛定理），并由单调收敛定理把
        不等式传递到极限。对变号或复值的 $f$，
        由 $\abs{Hf} \leq H\abs f$ 得出结论。
\end{enumerate}

\medskip
\noindent\textbf{第二部分 --- 最优性。}
\begin{enumerate}[resume]
  \item 对 $A > 1$，令 $f_A(t) = t^{-1/p}\,\mathbf
        1_{\intcc1A}(t)$。计算 $\norm{f_A}_p^p = \ln A$
        并对 $1 \leq x \leq A$，
        \[
        (Hf_A)(x) = \frac p{p-1}\;x^{-1/p}\,
        \bigl(1 - x^{-(1 - 1/p)}\bigr).
        \]
  \item 由此可得 $\liminf_{A\to\infty} \norm{Hf_A}_p/
        \norm{f_A}_p \geq \frac{p}{p-1}$，从而
        断定该常数是最优的。
  \item 证明等式 $\norm{Hf}_p = \frac
        p{p-1}\norm f_p$ 在 $f \neq 0$ 时不可能成立。
        \emph{（考察赫尔德不等式在第一部分第 3 问中的等号情形，
        这将迫使 $f = cx^{-1/p}$ 型，而此类
        函数不属于 $L^p$。）}
\end{enumerate}

\medskip
\noindent\textbf{第三部分 --- 离散不等式。}
\begin{enumerate}[resume]
  \item 设 $g \geq 0$ 是 $(0,\infty)$ 上的非增函数，
        令 $a_n = g(n)$，比较 $\sum a_n^p$ 与 $\int g^p$，
        并相应地比较 $H$ 的平均，由第一部分推出
        \emph{哈代离散不等式}：对 $a_n \geq 0$，
        \[
        \sum_{n\geq1}\Bigl(\frac{a_1 + \dots +
        a_n}{n}\Bigr)^{p}
        \;\leq\;
        \Bigl(\frac{p}{p-1}\Bigr)^{p}\,\sum_{n\geq1}a_n^p
        \]
        --- 先对非增数列 $(a_n)$ 利用上面的比较证明它，
        再通过重排把一般情形化为
        非增情形，\emph{（承认，附一行理由
        ：将 $(a_n)$ 按递减顺序排列
        只会使左端增大，
        同时保持右端不变。）}
  \item 由此可得：若 $\sum a_n^p < \infty$，则 $(a_n)$ 的切萨罗
        平均仍属于 $\ell^p$ --- 并给出一例
        （$p = 2$），其中 $(a_n) \in \ell^2$，而 $a_n$
        不可求和，但哈代不等式仍控制其平均。
\end{enumerate}

\medskip
\noindent\textbf{第四部分 --- 尾声。}
\begin{enumerate}[resume]
  \item 证明哈代不等式在 $p = 1$ 时不成立：取
        $f = \mathbf 1_{\intcc01}$，计算 $Hf$，并注意到
        $Hf \notin L^1$。证明在何处失效？
\end{enumerate}

\medskip
\noindent\textbf{第五部分 --- 极大函数，以及
勒贝格微分定理。} 哈代从原点取平均；
哈代--利特尔伍德则在\emph{每一点附近}取平均。
对 $f \in L^1(\R)$，定义
\[
Mf(x) = \sup_{r>0}\ \frac1{2r}\int_{x-r}^{x+r}\abs
f\,\dd\lambda .
\]
\begin{enumerate}[resume]
  \item （维塔利，有限情形）设 $B_1, \dots, B_N$ 为
        开区间。证明存在两两不交的子族
        $B_{i_1}, \dots, B_{i_k}$，使得 $\bigcup_jB_j
        \subseteq \bigcup_l3B_{i_l}$，其中 $3B$ 表示
        与原区间同心、长度为其三倍的区间，
        \emph{（贪心法：反复选取
        与已选区间不相交的最长区间）}。
  \item （弱型 $(1,1)$）证明：对每个 $t > 0$，
        \[
        \lambda\bigl(\{Mf > t\}\bigr) \;\leq\;
        \frac3t\,\norm f_1 :
        \]
        对每个 $x$，若 $Mf(x) > t$，则存在以之为中心的区间
        $B_x$，使 $\int_{B_x}\abs f > t\,\lambda(B_x)$；取
        紧集 $K \subseteq \{Mf > t\}$（内正则性），
        用有限个 $B_x$ 覆盖它，应用第 11 题，
        并穷竭之。
  \item 计算 $M\mathbf 1_{\intcc01}(x)$（$x > 1$），并推出
        $Mf \notin L^1$ 对每个 $f \neq 0$ 成立
        （$Mf(x) \geq \frac c{\abs x}$，于无穷远处）。当 $p =
        1$ 时，第 12 题的弱型不等式
        已是最优结论。
  \item （$p > 1$ 时的强型）对 $f \in L^p$，将 $f =
        f\,\mathbf 1_{\abs f > t/2} + f\,\mathbf 1_{\abs f
        \leq t/2}$，注意到 $Mf \leq M\bigl(f\mathbf 1_{\abs
        f > t/2}\bigr) + \frac t2$，并结合第 12 题的
        结论与层饼公式
        （见\cref{prop:b3:product:layercake}）及托内利定理，
        证明下式：
        \[
        \norm{Mf}_p^p \;\leq\;
        \frac{6p\,2^{p-1}}{p-1}\,\norm f_p^p .
        \]
        （当 $p \downarrow 1$ 时的爆破，即第 13 题中
        失效现象的定量化。）
  \item （勒贝格微分定理）证明：对 $f \in
        L^1(\R)$,
        \[
        \frac1{2r}\int_{x-r}^{x+r}f\,\dd\lambda
        \;\longrightarrow\; f(x)
        \qquad (r \to 0)\quad\text{对几乎处处的\ }x .
        \]
        \emph{（对连续的 $f$ 显然成立。一般地，记 $f
        = g + h$，$g$ 为具紧支集的连续函数，
        $\norm h_1 < \varepsilon$
        （见\cref{thm:b3:lp:density}）；平均振幅的
        $\limsup_{r\to0}$ 超过 $\delta$ 的点集
        位于 $\{Mh > \delta/2\} \cup \{\abs h
        > \delta/2\}$ 中，测度为 $O(\varepsilon/\delta)$；
        令 $\varepsilon \to 0$，再沿某一数列
        令 $\delta \to 0$。）}
  \item 由此推出：(a) 几乎每一点都是\emph{勒贝格
        点}（相对于 $f$）；(b) 对 $f \in L^1$，原函数
        $F(x) = \int_0^xf$ 几乎处处可导，且 $F' =
        f$ 几乎处处成立 --- 勒贝格意义下的微积分
        基本定理的积分部分，闭合了
        由阶梯函数开启的循环，
        （见\cref{pb:b3:measure:1}），表明其逆向部分
        可以不成立。
  \item （密度点）对可测集 $A \subseteq \R$，
        证明几乎每个 $x \in A$ 满足
        $\frac{\lambda(A\cap\intcc{x-r}{x+r})}{2r} \to 1$:
        可测集在其几乎所有的点处局部充满。
        用两行草拟如何由此得到
        施泰因豪斯定理的又一种证明
        （见\cref{exo:b3:measure:8}）。
\end{enumerate}

\medskip
\noindent\textbf{第六部分 --- 平均主题的
变奏。}
\begin{enumerate}[resume]
  \item （加权哈代不等式）对 $\alpha < p - 1$ 及 $f \geq 0$，
        证明下式：
        \[
        \int_0^\infty\Bigl(\frac{F(x)}x\Bigr)^{p}
        x^{\alpha}\,\dd x \;\leq\;
        \Bigl(\frac{p}{p - 1 - \alpha}\Bigr)^{p}
        \int_0^\infty f(x)^p\,x^{\alpha}\,\dd x
        \]
        用相同的分部积分法，并验证
        临界情形 $\alpha = p - 1$ 确实不能成立
        （改造第 10 题中的反例）。
  \item （伴随算子）记 $H^*f(x) =
        \int_x^{\infty}\frac{f(t)}t\,\dd t$。证明 $\langle
        Hf, g\rangle = \langle f, H^*g\rangle$，对
        非负的 $f, g$（由托内利定理），并证明
        $\norm{H^*f}_p \leq p\,\norm f_p$ \emph{（直接分部积分，
        或由共轭指数上的哈代不等式经
        对偶得到 --- 注意哪个指数携带哪个
        常数）}。
  \item （希尔伯特型不等式）由此推出，对
        非负的 $f \in L^p$，$g \in L^q$：
        \[
        \int_0^\infty\!\!\int_0^\infty
        \frac{f(x)\,g(y)}{\max(x,y)}\,\dd x\,\dd y
        \;\leq\; (p + q)\,\norm f_p\,\norm g_q
        \]
        \emph{（沿 $y < x$ 与 $y \geq x$ 分开讨论：每一半都是
        某一函数与另一函数的哈代变换的
        配对）}。
  \item （离散情形的最优性）证明第 8 题中的常数
        $\bigl(\frac p{p-1}\bigr)^p$ 同样是最优的，
        取 $a_n = n^{-1/p}\,\mathbf 1_{n \leq
        N}$，将两端与积分相比较，并令 $N
        \to \infty$（第二部分的离散镜像）。
  \item （综合）本题出现了三种平均算子：
        哈代算子 $H$、离散切萨罗
        平均与极大算子 $M$。各用一行说明其有界性的含义，
        注意到三者恰在 $p = 1$ 时全部失效，
        并解释为何这实为同一失效现象的三次
        重现（来自调和尾
        $\frac1x$）。
\end{enumerate}

\medskip
\noindent\textbf{第 VII 部分 --- 卡莱曼不等式，以及常数的
最优程度究竟如何。}
\begin{enumerate}[resume]
  \item （卡莱曼不等式）设 $a_n \geq 0$ 且 $\sum
        a_n < \infty$。对第 8 题的离散哈代不等式
        应用于 $b_n = a_n^{1/p}$，利用算术--几何平均不等式，并
        令 $p \to
        \infty$（证明 $p \mapsto
        \bigl(\frac p{p-1}\bigr)^p$ 递减至 $\eu$），由此
        得到
        \[
        \sum_{n\geq1}\bigl(a_1a_2\cdots a_n\bigr)^{1/n}
        \;\leq\; \eu\,\sum_{n\geq1}a_n :
        \]
        可和数列的几何平均仍
        可和，代价至多为 $\eu$。
  \item （常数 $\eu$ 是最优的）取 $a_n =
        \frac1n\,\mathbf 1_{n\leq N}$，利用斯特林
        夹逼（见\cref{pb:b3:product:1}），证明
        $(n!)^{-1/n} = \frac\eu n\bigl(1 +
        O\bigl(\frac{\ln n}n\bigr)\bigr)$，由此推出
        卡莱曼不等式两端均以 $\eu\ln N$ 的速度增长，并
        得出任何小于 $\eu$ 的常数都不成立。（注意
        规律是：哈代与卡莱曼的最优化序列
        都是恰好不属于该空间的调和型序列，
        即临界情形。）
  \item （“最优”的逼近有多慢？）取 $p = 2$。
        对 $f = \mathbf 1_{\intcc01}$，计算
        $\norm{Hf}_2/\norm f_2 = \sqrt2$，与界
        $2$ 相比。对第 5 题的近似最优元 $f_A$，
        证明如下精确等式
        \[
        \norm{Hf_A}_2^2 = 4\ln A - 8 + \frac{8}{\sqrt A},
        \qquad\text{故}\qquad
        \frac{\norm{Hf_A}_2^2}{\norm{f_A}_2^2}
        = 4 - \frac{8 - 8A^{-1/2}}{\ln A} .
        \]
        代入 $A = \eu^{10}$（比值 $\approx 1.790$）
        并评论：上确界 $2$ 只以
        $1/\ln A$ 的速度被逼近 --- 最优常数在数值上几乎
        无法察觉。
\end{enumerate}
\end{problem}
